\documentclass[a4paper,11pt]{article}

\newcommand{\TPnumero}{Lab 2 -- Solution}
% =========================================================
% Tutorial / lab preamble — Time Series Analysis module
% article class + fancyhdr + tcolorbox + listings (English)
% Author : Aymen Ben Brik (aymen.benbrik@esprit.tn)
% =========================================================

\usepackage[utf8]{inputenc}
\usepackage[T1]{fontenc}
\usepackage[english]{babel}
\usepackage{lmodern}
\usepackage{amsmath, amssymb, amsthm, mathtools, bm}
\usepackage{graphicx}
\usepackage{booktabs}
\usepackage{array}
\usepackage{multirow}
\usepackage{colortbl}
\usepackage{xcolor}
\usepackage{tikz}
\usetikzlibrary{positioning, arrows.meta, calc, shapes, fit, backgrounds, matrix}
\usepackage{listings}
\usepackage[most]{tcolorbox}
\usepackage{geometry}
\geometry{a4paper, margin=2cm, headheight=15pt}
\usepackage{fancyhdr}
\usepackage[hidelinks]{hyperref}
\usepackage{enumitem}

% --- Colours ---
\definecolor{espritBlue}{RGB}{30, 60, 110}
\definecolor{espritAccent}{RGB}{200, 80, 30}
\definecolor{boxEx}{RGB}{30, 60, 110}
\definecolor{boxQ}{RGB}{180, 120, 0}
\definecolor{boxC}{RGB}{30, 100, 60}
\definecolor{boxAtt}{RGB}{200, 80, 30}

% --- Listings ---
\definecolor{codebg}{RGB}{248,248,248}
\definecolor{codeframe}{RGB}{220,220,220}
\definecolor{keyword}{RGB}{0,82,155}
\definecolor{comment}{RGB}{88,110,117}
\definecolor{string}{RGB}{133,153,0}

\lstdefinestyle{pythonstyle}{
  language=Python,
  basicstyle=\ttfamily\small,
  keywordstyle=\color{keyword}\bfseries,
  commentstyle=\color{comment}\itshape,
  stringstyle=\color{string},
  backgroundcolor=\color{codebg},
  frame=single,
  rulecolor=\color{codeframe},
  frameround=tttt,
  breaklines=true,
  showstringspaces=false,
  tabsize=2
}
\lstdefinestyle{Rstyle}{
  language=R,
  basicstyle=\ttfamily\small,
  keywordstyle=\color{keyword}\bfseries,
  commentstyle=\color{comment}\itshape,
  stringstyle=\color{string},
  backgroundcolor=\color{codebg},
  frame=single,
  rulecolor=\color{codeframe},
  frameround=tttt,
  breaklines=true,
  showstringspaces=false,
  tabsize=2
}
\lstset{style=pythonstyle}

% --- Common macros (configurable path) ---
\providecommand{\macrosPath}{../../preamble/macros.tex}
\input{\macrosPath}

% --- Exercise counter ---
\newcounter{excounter}
\renewcommand{\theexcounter}{\arabic{excounter}}

\newtcolorbox[use counter=excounter]{exercise}[1][]{
  enhanced, breakable,
  colback=boxEx!5, colframe=boxEx, fonttitle=\bfseries,
  title={Exercise~\theexcounter\ifstrempty{#1}{}{ -- #1}},
  arc=2pt, boxrule=0.6pt
}
\newtcolorbox{question}[1][]{
  enhanced, breakable,
  colback=boxQ!5, colframe=boxQ, fonttitle=\bfseries,
  title={Question\ifstrempty{#1}{}{ -- #1}},
  arc=2pt, boxrule=0.5pt
}
\newtcolorbox{solution}[1][]{
  enhanced, breakable,
  colback=boxC!5, colframe=boxC, fonttitle=\bfseries,
  title={Solution\ifstrempty{#1}{}{ -- #1}},
  arc=2pt, boxrule=0.5pt
}
\newtcolorbox{warning}[1][]{
  enhanced, breakable,
  colback=boxAtt!5, colframe=boxAtt, fonttitle=\bfseries,
  title={Warning\ifstrempty{#1}{}{ -- #1}},
  arc=2pt, boxrule=0.5pt
}

% --- Headers / footers ---
\pagestyle{fancy}
\fancyhf{}
\fancyhead[L]{\textsc{\moduleNom} -- \auteurNom}
\fancyhead[R]{\TPnumero}
\fancyfoot[C]{\thepage}
\fancyfoot[L]{\auteurInstitution}
\fancyfoot[R]{\auteurEmail}
\renewcommand{\headrulewidth}{0.4pt}
\renewcommand{\footrulewidth}{0.2pt}

% --- Placeholder ---
\providecommand{\TPnumero}{Lab}


\title{Lab 2: Decomposition\\\large \emph{Reference solution}}
\author{\auteurNom\\\auteurEmail\\\auteurInstitution}
\date{\anneeUniversitaire}

\begin{document}
\maketitle

\noindent\emph{Reference Python code is in \texttt{correction.py}; this
document gives the analytical answers and expected numerical values.}

\setcounter{excounter}{0}
\begin{exercise}[Synthetic decomposition]\end{exercise}

\begin{solution}
The 4-panel figure shows the observed series with a clear linear drift
and a regular 12-month sinusoidal pattern of constant amplitude $5$ —
the visual signature of an \emph{additive} model.
\end{solution}

\begin{exercise}[Trend estimation: three methods]\end{exercise}

\begin{solution}
Reference RMSE values (over $T = 240$, seed 42):
\begin{center}
\begin{tabular}{lc}
\toprule
\textbf{Method}                                  & \textbf{RMSE vs.\ true trend} \\
\midrule
Centred MA, $k = 6$                              & $\approx 0.39$ \\
Linear regression (correctly specified)          & $\approx 0.07$ \\
LOESS, $\mathrm{frac} = 0.2$                     & $\approx 0.31$ \\
\bottomrule
\end{tabular}
\end{center}
The linear regression is the best estimator here because it matches
the true generating model. The MA loses a few observations at each
end and absorbs some sinusoidal residual at length 13. LOESS provides
flexibility but adds variance. The polynomial estimate is a clean
straight line; LOESS wiggles slightly around it.
\end{solution}

\begin{exercise}[Seasonality estimation: three methods]\end{exercise}

\begin{solution}
After LOESS detrending, the three methods produce visually similar
estimates of the 12-month seasonal pattern. RMSE versus the truth
$5 \sin(2\pi k / 12)$ (seed 42):
\begin{itemize}
  \item Seasonal averaging: $\approx 0.31$.
  \item Dummy regression (11 dummies): $\approx 0.31$ (mathematically
        identical to averaging when no other regressor is present).
  \item Cosine--sine model (1 harmonic): $\approx 0.16$ — \textbf{best},
        using only 2 parameters.
\end{itemize}
The Fourier model wins here because the truth is a \emph{pure}
sinusoid, exactly captured by the 1-harmonic basis (which lets the
estimator pool all 240 observations into 2 coefficients). Averaging
and dummies estimate 12 independent indices, with one degree of
freedom per index and therefore higher variance. With more harmonics
($J \geq 2$), the Fourier model can fit more complex seasonal shapes;
with all $J = \lfloor d/2 \rfloor$ harmonics it is equivalent to the
dummies (Proposition~2.5 of the lecture).
\end{solution}

\begin{exercise}[Application to Atlanta temperatures]\end{exercise}

\begin{solution}
\textbf{(1)} Period $d = 12$ (annual seasonality). The amplitude of
the seasonal swing looks roughly constant over the 138-year span,
supporting an additive model.

\textbf{(2)} Visual decomposition: the trend (centred MA, $k = 6$)
sits around $61.8^\circ$F. The 12 monthly seasonal indices, in
$^\circ$F, are
\[
\begin{array}{rrrrrrrrrrrr}
\text{Jan} & \text{Feb} & \text{Mar} & \text{Apr} & \text{May} & \text{Jun}
& \text{Jul} & \text{Aug} & \text{Sep} & \text{Oct} & \text{Nov} & \text{Dec}\\
-19.9 & -17.1 & -9.5 & -0.3 & +8.7 & +16.2 & +18.6 & +17.7 & +12.2 & +1.2 & -9.9 & -17.9
\end{array}
\]
Peak in July (+18.6), trough in January ($-19.9$), peak-to-peak
amplitude $\sim 38^\circ$F. The residual has std $\approx 2.9^\circ$F
— mostly weather noise.

\textbf{(3)} \texttt{statsmodels.seasonal\_decompose(y, period=12,
model="additive")} gives the same monthly seasonal indices (within
floating-point tolerance) when using the same $2 \times 12$ MA filter.
\end{solution}

\begin{exercise}[Bonus: additive vs multiplicative]\end{exercise}

\begin{solution}
The multiplicative simulated series shows a seasonal amplitude that
\emph{grows} with the trend level. The additive workflow on $Y_t$
leaves a clearly heteroscedastic residual (funnel-shaped), whereas
the same workflow on $\log Y_t$ gives an approximately stationary
residual.

\textbf{Conclusion:} prefer the multiplicative model (or work on
$\log Y_t$) when the seasonal amplitude scales with the level — typical
of sales, GDP, population, financial volume.
\end{solution}

\end{document}
